% ============================================================================
% 4.3 uart_tx.v
% ============================================================================
\subsection{设计目标}

接收一个并行的 8 位字节（\texttt{tx\_data}），在 \texttt{tx\_data\_valid = 1}
时捕获它，然后按 UART 协议逐位发送到 \texttt{tx\_pin} 上。

协议：

一帧 11 位（起始位 + 8 数据位 + 停止位 + 空闲）：

\begin{center}
\begin{tabular}{lllllllllll}
位 & 0 & 1--8 & 9 & 10 \\
作用 & 起始 & 数据 b0--b7 & 停止 & 空闲 \\
电平 & 低 & (数据) & 高 & 高 \\
\end{tabular}
\end{center}

每一位的持续时间由波特率决定。在 50 MHz 时钟、115200 波特率下，
一位持续 \( \frac{50,000,000}{115200} \approx 434 \) 个时钟周期。

\subsection{状态机设计}

发送器使用 3 个状态的简单状态机：

\begin{center}
\begin{tabular}{l l}
\textbf{S\_IDLE} & 空闲。等待 \texttt{tx\_data\_valid}。 \\
\textbf{S\_START} & 发送起始位（拉低 tx\_pin 一位的时间）。 \\
\textbf{S\_SEND\_BYTE} & 逐位发送 8 个数据位，低位在前。 \\
\textbf{S\_STOP} & 发送停止位（拉高 tx\_pin 一位的时间），完成后回到 IDLE。 \\
\end{tabular}
\end{center}

\subsection{源码与讲解}

\begin{lstlisting}[style=verilog,caption=Chip/src/uart\_tx.v]
module uart_tx
#(
    parameter CLK_FRE = 50,      // 时钟频率(MHz)
    parameter BAUD_RATE = 115200 // 波特率
)
(
    input clk,
    input rst_n,
    input [7:0] tx_data,         // 要发送的字节
    input tx_data_valid,         // 数据有效：一次握手的开始
    output reg tx_data_ready,    // 告诉发送方：我已准备好接收下一个
    output tx_pin                // 串行输出引脚
);

localparam CYCLE = CLK_FRE * 1000000 / BAUD_RATE;
// CYCLE = 434 (50 MHz / 115200)

// 状态编码
localparam S_IDLE      = 1;
localparam S_START     = 2;
localparam S_SEND_BYTE = 3;
localparam S_STOP      = 4;

reg [2:0]  state;
reg [15:0] cycle_cnt;    // 每一位内的时钟计数（最多 65535）
reg [2:0]  bit_cnt;      // 已经发送了第几位（0..7）
reg [7:0]  tx_data_latch; // 锁存住要发送的字节
reg        tx_reg;        // 输出寄存器

assign tx_pin = tx_reg;   // tx_reg 的值直接送到 tx 引脚
\end{lstlisting}

状态转移逻辑：

\begin{lstlisting}[style=verilog]
always @(posedge clk or negedge rst_n)
begin
    if (!rst_n)
        state <= S_IDLE;
    else
        case (state)
            S_IDLE:
                if (tx_data_valid)
                    state <= S_START;
                else
                    state <= S_IDLE;
            S_START:
                if (cycle_cnt == CYCLE - 1)
                    state <= S_SEND_BYTE;
                else
                    state <= S_START;
            S_SEND_BYTE:
                if (cycle_cnt == CYCLE - 1 && bit_cnt == 3'd7)
                    state <= S_STOP;
                else
                    state <= S_SEND_BYTE;
            S_STOP:
                if (cycle_cnt == CYCLE - 1)
                    state <= S_IDLE;
                else
                    state <= S_STOP;
        endcase
end
\end{lstlisting}

数据准备好标志和锁存：

\begin{lstlisting}[style=verilog]
always @(posedge clk or negedge rst_n)
begin
    if (!rst_n)
        tx_data_ready <= 1'b0;
    else if (state == S_IDLE)
        tx_data_ready <= !(tx_data_valid);  // 空闲时才能接收
    else if (state == S_STOP && cycle_cnt == CYCLE - 1)
        tx_data_ready <= 1'b1;
end

always @(posedge clk or negedge rst_n)
begin
    if (!rst_n)
        tx_data_latch <= 8'd0;
    else if (state == S_IDLE && tx_data_valid)
        tx_data_latch <= tx_data;    // 在进入 START 状态时锁存数据
end
\end{lstlisting}

计数器（计算当前位已发送了多少个时钟周期）：

\begin{lstlisting}[style=verilog]
always @(posedge clk or negedge rst_n)
begin
    if (!rst_n)
        cycle_cnt <= 16'd0;
    else if ((state == S_START || state == S_SEND_BYTE ||
              state == S_STOP) && cycle_cnt < CYCLE - 1)
        cycle_cnt <= cycle_cnt + 16'd1;
    else
        cycle_cnt <= 16'd0;
end

always @(posedge clk or negedge rst_n)
begin
    if (!rst_n)
        bit_cnt <= 3'd0;
    else if (state == S_SEND_BYTE)
        if (cycle_cnt == CYCLE - 1)
            bit_cnt <= bit_cnt + 3'd1;
        else
            bit_cnt <= bit_cnt;
    else
        bit_cnt <= 3'd0;
end
\end{lstlisting}

最后，真正驱动输出引脚的逻辑：

\begin{lstlisting}[style=verilog]
always @(posedge clk or negedge rst_n)
begin
    if (!rst_n)
        tx_reg <= 1'b1;   // UART 空闲时 tx_pin 为高电平
    else
        case (state)
            S_IDLE:      tx_reg <= 1'b1;   // 空闲=高
            S_START:    tx_reg <= 1'b0;   // 起始位=低
            S_SEND_BYTE:
                         tx_reg <= tx_data_latch[bit_cnt];  // 从 bit0 开始逐位发
            S_STOP:     tx_reg <= 1'b1;   // 停止位=高
            default:    tx_reg <= 1'b1;
        endcase
end
endmodule
\end{lstlisting}

\paragraph{锁存为什么重要}
\texttt{tx\_data\_latch} 在进入 S\_START 时把 \texttt{tx\_data} 的值存下来。
如果没有这一步，上游逻辑在发送过程中改变 \texttt{tx\_data}，
我们就会输出一个被污染的字节——先几位是旧值，后几位是新值。

\paragraph{为什么 UART 空闲时 tx\_pin 必须为高}
UART 协议里，"线路空闲"就是高电平。起始位是把线路拉低一位时间，
这样接收端才能检测到"有数据开始到来"。如果空闲时电平是低，
接收端会把空闲状态误当做起始位，整个通信就错乱了。
